\notechapter{N-20230220}{Special Proofs of Trigonometric Addition Formulas}

\section{Motivation}

This handwritten note tries several proofs of trigonometric addition formulas.  The author first tried a mixed line-slope proof for tangent, then tried reciprocal bases, then settled on more reliable methods for sine and cosine.  This history is important: it records why some apparently clever approaches become awkward.

\section{Sine difference formula by geometry}

Place two unit vectors at angles \(\alpha\) and \(\beta\):
\[
  u=(\cos\alpha,\sin\alpha),\qquad v=(\cos\beta,\sin\beta).
\]
The magnitude of the two-dimensional cross product is
\[
  |u\times v|=\sin|\beta-\alpha|.
\]
But the determinant gives
\[
  u\times v=\cos\alpha\sin\beta-\sin\alpha\cos\beta.
\]
Hence, with orientation taken into account,
\[
  \sin(\beta-\alpha)=\sin\beta\cos\alpha-\cos\beta\sin\alpha.
\]

This proof is clean because the determinant already measures signed area, and signed area of two unit vectors is sine of the included angle.

\section{Rotation-matrix proof}

The rotation matrix is
\[
  R(\theta)=
  \begin{pmatrix}
  \cos\theta&-\sin\theta\\
  \sin\theta&\cos\theta
  \end{pmatrix}.
\]
Because rotating by \(\alpha\) and then by \(\beta\) is the same as rotating by \(\alpha+\beta\),
\[
  R(\alpha+\beta)=R(\beta)R(\alpha).
\]
Multiplying matrices gives
\[
  \cos(\alpha+\beta)=\cos\alpha\cos\beta-\sin\alpha\sin\beta,
\]
\[
  \sin(\alpha+\beta)=\sin\alpha\cos\beta+\cos\alpha\sin\beta.
\]
This is probably the most conceptual elementary proof: addition formulas are simply composition laws for rotations.

\section{Euler formula proof}

Euler's formula gives
\[
  e^{i\theta}=\cos\theta+i\sin\theta.
\]
Then
\[
  e^{i(\alpha+\beta)}=e^{i\alpha}e^{i\beta}.
\]
Expanding the right-hand side gives
\[
  (\cos\alpha+i\sin\alpha)(\cos\beta+i\sin\beta)
\]
\[
  =(\cos\alpha\cos\beta-\sin\alpha\sin\beta)
  +i(\sin\alpha\cos\beta+\cos\alpha\sin\beta).
\]
Equating real and imaginary parts yields the addition formulas.

\section{Tangent formula}

Once sine and cosine are known,
\[
  \tan(\alpha+\beta)=\frac{\sin(\alpha+\beta)}{\cos(\alpha+\beta)}
  =\frac{\tan\alpha+	an\beta}{1-	an\alpha\tan\beta},
\]
provided the denominators are nonzero.  This is why trying to prove tangent first is often less natural: tangent is a quotient, while sine and cosine are coordinates of rotation.

\section{The method behind the trick}

The same formulas have three interpretations: geometry says sine is signed area, linear algebra says rotations compose by matrix multiplication, and complex analysis says rotations are multiplication by unit complex numbers.  The formulas are memorable once they are seen as structure-preservation laws rather than isolated identities.

\section{Rotation proof as the structural proof}

The cleanest proof of trigonometric addition formulas is the identity
\[
  R(a+b)=R(a)R(b)
\]
for plane rotations.  Matrix multiplication gives
\[
  \cos(a+b)=\cos a\cos b-\sin a\sin b,
\]
\[
  \sin(a+b)=\sin a\cos b+\cos a\sin b.
\]
This proof shows that the formulas express composition of rotations.

\section{Euler formula and characters}

Euler's formula
\[
  e^{i(a+b)}=e^{ia}e^{ib}
\]
gives the same identities by comparing real and imaginary parts.  The map \(t\mapsto e^{it}\) is a character from the additive group of real angles to the multiplicative circle group.

\section{Tangent formula as projective coordinate}

Tangent is the slope coordinate on the unit circle, where defined.  The formula
\[
  \tan(a+b)=\frac{\tan a+\tan b}{1-\tan a\tan b}
\]
is the induced fractional-linear rule for slopes under rotation.

\section{Addition formulas as rotation composition}

The many proofs of addition formulas are different coordinate expressions of one fact: rotations form a group.  Geometry, matrices, complex exponentials, and tangent slopes all encode the same composition law.
